\documentclass{ximera}[font=12pt]  % I don't think font sizing works here

\title{Fractions: GCF}   % Use xmlatex all -s to compile when SERVE refuses to work.
\author{Kelly Stady}
\license{CC: 0}         % replace with an appropriate license, or set it in xmPreamble

\begin{document}


\begin{abstract}

\end{abstract}

\maketitle

\label{xim:FracsGCFSimplify}


\section*{Greatest Common Factor (GCF)}

In the next section, we will simplify fractions using the largest number that evenly divides into both the numerator and denominator. 

\begin{definition}
The \textbf{greatest common factor} is the largest whole number that divides evenly into two or more numbers.  Equivalently, it is
the largest whole number that is a factor of two or more numbers.

\end{definition}


\begin{problem} \textcolor{orange!80!black}{\textbf{True or False}}  The GCF will always be less than or equal to the smallest 
    number in the group.

\begin{multipleChoice}
\choice[correct]{True}
\choice{False}
\end{multipleChoice}

\begin{feedback}[correct]
That's right!  Since the GCF evenly divides into each number, it can't be larger than the smallest number.
\end{feedback}

\end{problem}


\begin{procedure}[\textbf{Finding the GCF Using Lists of Factors}]

List the factors of each number and look for the largest number that appears in each list.  \textbf{Alternatively}, 
only list the factors of the smallest number then look for the largest factor that evenly divides into the other number(s).

\end{procedure}


\begin{example}

Find the GCF of 12 and 36. 

\begin{explanation}

We begin by listing the factors of each number.

\begin{flushleft}
\hspace*{8em} % Move this OUTSIDE the math block for better control
$\begin{aligned}
& \textbf{Factors of 12}: \ 1, 2, 3, 4, 6, 12 \\[1.5ex]
& \textbf{Factors of 36}: \ 1, 3, 4, 6, 9, 12, 18, 36
\end{aligned}$
\end{flushleft}

The \textit{common factors} of 12 and 36 are 1, 3, 4, 6, and 12.  So, the \textbf{greatest common factor} is 12.  


\begin{remark}
Notice that it is more efficient to list the factors of 12 and look for the largest number that evenly divides into 36.  
\end{remark}

\end{explanation}

\end{example}


\begin{problem}
Find the GCF of $20$ and $30.$ 

\begin{enumerate}
\item List the factors of $20$ and $30.$

\begin{multipleChoice}
    \choice{\textbf{Factors of 20}: 1, 2, 4, 5, 20; \ \textbf{Factors of 30}: 1, 2, 3, 5, 6, 10, 15, 30}
    \choice{\textbf{Factors of 20}: 1, 2, 4, 5, 10, 20; \ \textbf{Factors of 30}: 1, 2, 3, 5, 10, 15, 30}
    \choice[correct]{\textbf{Factors of 20}: 1, 2, 4, 5, 10, 20; \ \textbf{Factors of 30}: 1, 2, 3, 5, 6, 10, 15, 30}
    \choice{\textbf{Factors of 20}: 2, 4, 5, 10; \ \textbf{Factors of 30}: 2, 3, 5, 6, 10, 15}
\end{multipleChoice}

\item Based on the lists above, what is the greatest common factor (GCF) of $20$ and $30$?

\[ \text{GCF} = \answer{10} \]

\end{enumerate}

\end{problem}



INSERT VIDEO SHOWING Alternative IDEA



\begin{problem}
Find the GCF of 15 and 20 by listing the factors of both numbers or just the number 15.  Use whichever technique you prefer.

\[ \textrm{GCF} = \answer{5} \]

\end{problem}


\begin{problem}
Find the GCF of 24 and 60 by listing the factors of both numbers or just the number 24.  Use whichever technique you prefer.

\[ \textrm{GCF} = \answer{12} \]

\end{problem}


\begin{procedure}[\textbf{Finding the GCF Using Prime Factorizations}]
    
    Find the prime factorization of each number, then multiply the common prime factors.  
    
\end{procedure}

\begin{example}

Find the GCF of 12 and 36 using prime factorizations. 

\begin{explanation}

We begin by finding the prime factorization of each number.

\begin{eqnarray*}
12 &=& 2 \cdot 2 \cdot 2 \cdot 3 \\
   & & \\
36 &=& 2 \cdot 2 \cdot 3 \cdot 3 \\
\end{eqnarray*}

The common factors are 2 and 3.  There are two factors of 2 and one factor of 3 in \textit{both} factorizations, therefore

\begin{eqnarray*}
\textrm{GCF} \  &=& 2 \cdot 2 \cdot 3 \\
   & & \\
   &=& \answer{12} \\
\end{eqnarray*}


\end{explanation}
\end{example}

\begin{problem}
Find the GCF of 60 and 75 using their prime factorizations given below.

\begin{eqnarray*}
60 &=& 2 \cdot 2 \cdot 3 \cdot 5 \\
   & & \\
75 &=& 3 \cdot 5 \cdot 3 \\
\end{eqnarray*}

GCF = $\answer{15}$ 

\begin{feedback}
Excellent!  The common factors are 3 and 5.  There is one factor of 3 and one factor of 5 in \textbf{both} factorizations, so the GCF = $3 \cdot 5$ which equals 15.
\end{feedback}

\end{problem}


\begin{problem}
Find the GCF of 85 and 150 using their prime factorizations shown below.


\begin{eqnarray*}
85 &=& 5 \cdot 17 \\
   & & \\
150 &=& 2 \cdot 3 \cdot 5 \cdot 5 \\
\end{eqnarray*}


GCF = $\answer{5}$

\end{problem}


INSERT VIDEO:  Finding the GCF of 3 numbers!!!!!

The GCF can be used to solve problems involving groups of things that need to be divided efficiently, with no waste.  Use what you've learned
about finding the GCF to solve the following application problems.

\begin{problem}
  You are creating identical goodie bags for a child's birthday party.  You have \textbf{18} Airheads, \textbf{24} Tootsie Rolls and
  \textbf{36} Kit-Kats.  What is the \textbf{greatest} number of identical goodie bags you can make if you want to use all of the candy?

  \[ \textrm{Greatest Number of Goodie Bags} = \answer{6} \] 

  \begin{hint}
    To solve this, you need to find the Greatest Common Factor (GCF) of 18, 24, and 36.  List the factors for each number and find the 
    largest one that appears in all three lists.
  \end{hint}

  Now that you know you can make identical goodie bags with no waste, how many of each candy will be in \textbf{each} bag?

  \begin{itemize}
      \item Airheads: $\answer{3}$
      \item Tootsie Rolls: $\answer{4}$
      \item Kit-Kats: $\answer{6}$
  \end{itemize}
  
  \begin{hint}
    Divide the total number of each candy by the number of bags (6). For example, $18 \div 6 = 3$ Airheads.
  \end{hint}


\end{problem}


\begin{problem}
  A carpenter has three wooden boards of different lengths: \textbf{8} ft, \textbf{12} ft and
      \textbf{20} ft.  The carpenter wants to cut all three boards into smaller pieces of equal length,
       so that there is no scrap wood left over. What is the longest possible length for these pieces?

  \[ \textrm{Longest Length } = \answer{4} \ \textrm{feet} \] 

  \begin{hint}
    Find the Greatest Common Factor (GCF) of 8, 12, and 20. This will be the largest number that divides evenly into all three lengths.
  \end{hint}

How many pieces will the carpenter have after cutting all three boards?

  \[ \textrm{Total Pieces } = \ \answer{10} \] 

  \begin{hint}
    Divide each board's length by 4 to see how many pieces come from each:
    \begin{itemize}
        \item $8 \div 4 = 2$ pieces
        \item $12 \div 4 = 3$ pieces
        \item $20 \div 4 = 5$ pieces
    \end{itemize}
    Now, add them.
  \end{hint}
\end{problem}



\end{document}